% Ad Familiares 10.6 (ad-familiares-10-06)
% Source: https://raw.githubusercontent.com/PerseusDL/canonical-latinLit/master/data/phi0474/phi056/phi0474.phi056.perseus-lat1.xml
% Fetched by scripts/fetch_latin.py

% Dateline (Perseus): Scr. Romae xiii K. Apr. vesperia. 711 (43).

\ciceroLetterOpener{CICERO PLANCO.}

\ciceroSection{1} quae locutus est Furnius noster de animo tuo in rem p., ea gratissima fuerunt senatui, p. R. probatissima ; quae autem tuae recitatae litterae sunt in senatu, nequaquam consentire cum Furni oratione visae sunt. pacis enim auctor eras; cum conlega tuus, vir clarissimus, a foedissimis, latronibus obsideretur, qui aut positis armis pacem petere debent aut, si pugnantes eam postulant, victoria pax non pactione parienda est. sed de pace litterae vel Lepidi vel tuae quam in partem acceptae sint, ex viro optimo, fratre tuo, et ex C. Furnio poteris cognoscere.

\ciceroSection{2} me autem impulit tui caritas ut, quamquam nec tibi ipsi consilium deesset, et fratris Furnique benevolentia fidelisque prudentia tibi praesto esset futura, vellem tamen meae quoque auctoritatis pro plurimis nostris necessitudinibus praeceptum ad te aliquod pervenire. crede igitur mihi, Plance, omnis, quos adhuc gradus dignitatis consecutus sis (es autem adeptus amplissimos) eos honorum vocabula habituros, non dignitatis insignia, nisi te cum libertate populi R. et cum senatus auctoritate coniunxeris. seiunge te, quaeso, aliquando ab iis, cum quibus te non tuum iudicium sed temporum vincla coniunxerunt.

\ciceroSection{3} complures in perturbatione rei p. consules dicti, quorum nemo consularis habitus est, nisi qui animo exstitit in rem p. consularis. talem igitur te esse oportet, qui primum te ab impiorum civium tui dissimillimorum societate seiungas, deinde te senatui bonisque omnibus auctorem, principem, ducem praebeas, postremo ut pacem esse iudices non in armis positis sed in abiecto armorum et servitutis metu. haec si et ages et senties, tum eris non modo consul et consularis, sed magnus etiam consul et consularis ; sin aliter, tum in istis amplissimis nominibus honorum non modo dignitas nulla erit sed erit summa deformitas. haec impulsus benevolentia scripsi paulo severius ; quae tu in experiendo ea ratione, quae te digna est, vera esse cognosces. D. xiii K. Apr.
